\chapter*{Introducción}
\addcontentsline{toc}{chapter}{Introducción}

\lettrine[lines=2]{L}{a} visión es el más noble y elevado de los sentidos que nos han sido otorgados, sin embargo, está limitada en cuanto puede discernir. Sabemos por experiencia que nos es imposible resolver objetos demasiado alejados, como los planetas, las estrellas, como tambien los objetos diminutos, como lo pueden ser los tejidos y las células. Sin embargo, gracias al conocimiento de la naturaleza de la luz hemos podido sortear esta dificultad en la capacidad de dicernir, por medio de la manipulación de la luz que emana de estos cuerpos — ya sea por emisión directa o por reflexión — de una manera tal que el objeto en un mayor detalle se manifieste en los dominios accesibles a la visión humana, permitiendonos dicernir detalles que sin instrumento nos serían ocultos. 

Llamamos imagén a la luz que, emanada desde el objeto, se proyecta sobre la superficie o dominio de observación. Una imagén no siempre es una representación del objeto, pero en los \emph{sistemas formadores de imagén} se procura que lo sea. Cuando a cada punto del objeto le correponde un punto en la imagén, y la forma del objeto y la forma de la imagen son similares, se dice que la imagén es perfecta (\cite{bornwolf1999} sección \S 5.2). 


El problema de la formación de la imagen perfecta, es de una gran importancia, tanto para dicernir los diferentes objetos de la naturaleza, como para sistemas fotograficos, de impresión, y de lítografia y mucho más aún en sistemas que requieren de una alta resolución como en el caso de la litografía, microscopia, y astronomía\dots


El problema de la formación perfecta de un objeto extendido, es un problema de natureleza general, que requiere que el problema de la imagén perfecta de un punto sea resuelto, debido a que como dijimos, la imagen perfecta de un objeto se da cuando cada punto del objeto $Q$ tiene como imagén un punto $Q'$ en el espacio imagén, y cuando se da esto decimos que el punto $Q$ tiene una imagén estigmática, que es $Q'$. Es por ello que por \emph{sistemas estigmáticos} son de gran interés en óptica, por \emph{sistema estigmático} entendemos nos refererimos a los sistemas que aseguran que un par de puntos sean conjugados, es decir que uno sea la imagén estigmática del otro. 


Esta definición del sistema estigmático, y la noción misma de ``punto'' 
sobre la que descansa, pertenece propiamente a la óptica geométrica, pues en
óptica ondulatoria un punto objeto no tiene como imagén un punto en el sentido matemático, sino que su imagen es una distribución extendida que llamamos \emph{PSF (Point Spread Function)} por la difracción.
Aún así, los sistemas que son estigmáticos en el regimén geométrico no dejan de ser importantes en el marco de la \emph{Teoría Ondulatoría de la Luz} debido a que las aberraciones geométricas tienen un gran impacto en la teoría ondulatoria de la formación de imágenes \cite{bornwolf1999},\cite{zernike1934}, pues estas se manifiestan como un una aberración de fase (camino óptico) sobre la pupila del sistema.
Dicha aberración modifica la función de dispersión puntual (PSF): en lugar de la
PSF ideal de un sistema limitado únicamente por difracción, el sistema produce
una PSF aberrada por los defectos en a geométria, por lo que es en general más extensa, y con menor
concentración de energía en el máximo principal \cite{}. La
corrección de las aberraciones geométricas determina cuán cerca puede hallarse
la imagén del límite impuesto por la sola difracción. Por ello los
sistemas estigmáticos conservan su interés aun en el marco ondulatorio, pues estan libres de aberración esférica.

Con esto justificamos el estudio de los sistemas estigmáticos en la óptica ondulatoria, y nos dedicamos a estudiar el más simple de ellos, el formado por dos 
medios con diferente indice de refracción separados por una superficie que aseugra el esigmátismo para un par de puntos conjugados $A$ y $A'$. Estas superficies son conocidas desde 1668 con el apendice la DIOVTRIQUE del \emph{Discours de la méthode} \cite{dioptrique}. Si bien, bajo la aproximación paraxial se usan parabolas y estas aseguran el esigmatismo para un par de puntos igualmente, esto debido a que las superificies partesianas en el regimen paraxial son parabolas, hecho del que se puede deducir la posición de los puntos estigmáticos para este tipo de sistemas.



El régimen paraxial, donde las superficies cartesianas se reducen a parábolas, corresponde también al régimen de baja apertura numérica, en el que la descripción escalar del campo es suficiente. Al aumentar la apertura numérica —condición indispensable para mejorar la resolución espacial—, las superficies cartesianas se alejan apreciablemente de las parábolas y, con ello, la dirección local de propagación varía de manera significativa sobre la pupila del sistema. En este régimen, la refracción en la interfaz dióptrica no actúa como una simple modulación escalar de amplitud: en cada punto de la superficie aparecen direcciones locales de polarización $(s)$ y $(p)$, coeficientes de Fresnel distintos $t_s \neq t_p$, y una transformación vectorial que depende de la posición. Como consecuencia, los diferentes estados de polarización pueden acoplarse entre sí, y componentes del campo inicialmente ausentes pueden emerger tras la transmisión~\cite{Quabis2000,Dorn2003}.

A las desviaciones entre la distribución de intensidad predicha por la teoría escalar y la obtenida mediante una descripción electromagnética completa se las denomina \emph{aberraciones vectoriales}~\cite{Shen2022,Ma2025}. Estas aberraciones son de especial relevancia en sistemas de alta apertura numérica como los empleados en litografía óptica, microscopía confocal y detección astronómica, donde se exige llevar la resolución al límite impuesto por la difracción. En esos contextos, la corrección de las aberraciones geométricas mediante superficies estigmáticas elimina la aberración de fase sobre la pupila, pero no garantiza una PSF ideal: las aberraciones vectoriales constituyen un mecanismo de degradación adicional, de origen puramente electromagnético.

El objetivo de este trabajo es demostrar que en un sistema libre de aberraciones geométricas —formado por una única superficie dióptrica estigmática—, existen aberraciones de origen vectorial que modifican la distribución de intensidad en el plano focal y pueden deteriorar la resolución efectiva del sistema. Para ello se calcula el campo electromagnético resultante de la transmisión, a través del dioptrio, del frente de onda esférico emitido desde el punto objeto del sistema, incorporando de manera explícita los efectos vectoriales de la refracción, y se compara con la predicción de la teoría escalar.

El trabajo se organiza como sigue. En el capítulo~\ref{ch:geometrica} se establece el marco geométrico del problema: se deriva la forma de las superficies cartesianas a partir de la condición de igual camino óptico y se obtienen las cantidades locales necesarias para describir la refracción. El capítulo~\ref{ch:ondulatoria} introduce la descripción electromagnética del campo, las condiciones de frontera y los coeficientes de Fresnel que determinan la transformación vectorial en la superficie. En el capítulo~\ref{ch:focal} se combinan estos elementos para construir la función pupila vectorial del sistema y propagar el campo hasta el plano focal, identificando las diferencias respecto a la predicción escalar. Finalmente, el capítulo~\ref{ch:esfera} aplica el formalismo al dioptrio esférico, cuyo caso límite de estigmatismo (puntos de Young–Weierstrass) permite validar los resultados y examinar en qué medida la geometría de la superficie incide en las aberraciones vectoriales observadas.

